\section{Method}
\label{sec:method}

\subsection{Problem Reformulation: Verifiable Program Synthesis}

We define the task as follows. Given a target CAD part represented by a ground-truth STEP file $G$ and an optional structured input $J$ (a Fusion360 reconstruction JSON or a synthetic parameter spec), produce a CadQuery program $P$ such that $\texttt{execute}(P)$ yields a STEP file $\hat{G}$ with $\text{IoU}(\hat{G}, G) \geq \tau$, where $\tau = 0.99$. The generated program $P$ is accepted as a verified pair only if this criterion is met.

This reformulation has two important consequences. First, acceptance is fully automated and requires no human annotation: the OCCT kernel computes the IoU exactly. Second, the acceptance threshold $\tau$ is strict enough to reject most common failure modes, including wrong extrusion direction, sign-flipped coordinate axes, and arc midpoints that are a few millimeters off their intended circle.

\begin{figure*}[t]
  \centering
  \includegraphics[width=\linewidth]{fig2_pipeline.png}
  \caption{\textbf{Verification-Guided CAD Data Curation Pipeline.}
  Main flow (left to right): structured input (GT STEP + JSON) $\to$ LLM generation skill $\to$ subprocess execution $\to$ OCCT IoU verifier $\to$ pass/fail decision $\to$ verified store. Failure paths: execution failure routes to the provider cascade (codex $\to$ openai $\to$ glm); near-misses (IoU $\in [0.97, 0.99)$) route to the repair or manual check skill and re-enter verification. Accepted pairs are stored as JSONL records with full provenance metadata.}
  \label{fig:pipeline}
\end{figure*}

\subsection{Verification-Guided Curation Pipeline}

The main pipeline is shown in Figure~\ref{fig:pipeline}. For each target part, the pipeline proceeds as follows.

\noindent\textbf{Generation.}
The LLM receives a prompt containing the structured input JSON (trimmed to fit context) and a system prompt instructing it to produce executable CadQuery code. We use a cascade of providers in priority order: Codex CLI (\texttt{gpt-5.3-codex}), OpenAI API (\texttt{gpt-5.2}), and GLM (\texttt{glm-4.6v}). The cascade advances to the next provider if the current one fails execution or produces IoU below $\tau$.

\noindent\textbf{Execution.}
The generated Python code is run in a sandboxed subprocess with \texttt{LD\_LIBRARY\_PATH} set to expose the OCCT shared libraries. A monkey-patch (\texttt{\_OCP\_HASHCODE\_FIX}) is injected to work around the CadQuery \texttt{.union()} hash collision bug. The last line of every generated program must be \texttt{result.val().exportStep("output.step")}. Execution timeout is 30 seconds.

\noindent\textbf{Verification.}
IoU is computed by importing both the ground-truth STEP and the generated STEP via \texttt{cq.importers.importStep}, extracting the solid with \texttt{.val()}, and computing:
\begin{equation}
\text{IoU} = \frac{V(G \cap \hat{G})}{V(G) + V(\hat{G}) - V(G \cap \hat{G})}
\end{equation}
where $V(\cdot)$ denotes volume and $\cap$ is computed via OCCT's \texttt{.intersect()} method. Visual verification (4-view HLR render plus human or LLM comparison) is performed on a subset and stored but not used as the acceptance gate in current runs.

\noindent\textbf{Decision.}
If IoU $\geq 0.99$, the pair is accepted. If $0.97 \leq \text{IoU} < 0.99$ (near-miss), it enters the repair queue. Otherwise it is discarded. Near-misses are periodically processed through the repair skill and manual checking skill.

\noindent\textbf{Record.}
Accepted pairs are written to a canonical JSONL record linking \texttt{raw\_step\_path}, \texttt{ops\_json\_path}, \texttt{gen\_step\_path}, \texttt{cq\_code\_path}, \texttt{views\_raw\_dir}, \texttt{views\_gen\_dir}, \texttt{iou}, \texttt{json\_type}, \texttt{source}, \texttt{timestamp}, and \texttt{complexity\_class}.

\subsection{Input JSON Formats}

Two structured input formats feed the generation skill.

\noindent\textbf{Fusion360 reconstruction JSON} encodes the full construction history exported from Fusion360: a timeline of entities including \texttt{Sketch} (with profiles, loops, arc/line curves parameterized in local sketch coordinates) and \texttt{ExtrudeFeature} (with extent type, distance, direction, and operation type such as \texttt{NewBody} or \texttt{Join}). Coordinates are in centimeters; our prompts instruct the model to multiply by 10 for CadQuery's millimeter units. This format is available for the full Fusion360 Gallery dataset~\cite{fusion360gallery}.

\noindent\textbf{Synthetic parameter JSON} is a flat key-value spec generated by our own primitive generators for operations such as revolve, shell, loft, polar array, rectangular array, chamfer, fillet, and cylinder:
\begin{verbatim}
{"op": "revolve_ring", "inner_r_mm": 28.5,
 "outer_r_mm": 32.9, "height_mm": 11.2}
\end{verbatim}
These are simpler to generate from than the Fusion360 JSON but cover a narrower set of shapes. The synthetic generators produce ground-truth STEP files programmatically, ensuring the IoU oracle is exact.

\subsection{Skill-Based Modular Design}

Rather than a single monolithic script, the pipeline is composed of reusable skills (Figure~\ref{fig:skills}). Each skill has a well-defined input/output contract, can be invoked independently, and fails gracefully without corrupting pipeline state.

\begin{table}[h]
\centering
\small
\caption{Skills and their contracts.}
\label{tab:skills}
\begin{tabular}{llll}
\toprule
\textbf{Skill} & \textbf{Input} & \textbf{Output} & \textbf{Failure} \\
\midrule
Generation & JSON + model & CadQuery \texttt{.py} & LLM refusal \\
Execution & \texttt{.py} path & STEP file & timeout, empty solid \\
Verification & STEP + GT STEP & IoU float & OCCT crash \\
Repair & \texttt{.py} + tag & Patched \texttt{.py} & pattern unknown \\
Manual check & \texttt{.py} + GT & Fixed \texttt{.py} & human judgment \\
Assembly & Verified JSONL & SFT JSONL & path errors \\
\bottomrule
\end{tabular}
\end{table}

The key advantage of this decomposition is handling \textbf{heterogeneous failure modes}.
Across our dataset, common failure patterns include: wrong extrusion plane sign (e.g., \texttt{extrude(d)} instead of \texttt{extrude(-d)} for XZ-plane parts); \texttt{radiusArc} direction ambiguity for non-quarter arcs; coordinate axis swaps in YZ-plane sketches; inner profile loops (holes) missed by the LLM; and the \texttt{NewBody} semantics of Fusion360---where the ground-truth STEP records only the most recently created body, not the accumulated union.
Each failure type has a different root cause and requires a different targeted fix.
A rigid workflow that applies a single repair strategy would handle at most one of these; the skill framework allows the right skill to be routed to each case.

\subsection{Arc Repair: A Detailed Example}

A recurring failure in complex Fusion360 parts is arc midpoint error.
CadQuery's \texttt{radiusArc(endpoint, radius)} fits a circular arc to two endpoints at a given radius, but the choice of which of the two possible arcs (major vs.\ minor, left vs.\ right) is determined by the sign of the radius---and this sign is frequently wrong.
For arcs that are neither quarter-circles nor semicircles, even a correct-sign \texttt{radiusArc} may pick a slightly different circle than the ground-truth arc due to floating-point sensitivity, producing IoU $\approx 0.988$--$0.995$ (a near-miss).

The correct fix uses \texttt{threePointArc(midpoint, endpoint)}, which uniquely determines the arc through three points. The midpoint is computed from the Fusion360 JSON arc parameters:
\begin{equation}
\text{mid} = c + r \cdot R(\hat{v},\, \Delta\theta/2)
\end{equation}
where $c$ is the arc center, $r$ is the radius, $\hat{v}$ is the reference vector (from center to start), $\Delta\theta$ is the sweep angle, and $R(\hat{v}, \alpha)$ rotates $\hat{v}$ by $\alpha$ in the sketch plane. For CCW arcs, $\Delta\theta > 0$; for CW traversal (when the profile's end point matches the arc's start), the arc is traversed in reverse and the midpoint formula uses $-\Delta\theta/2$. This computation is exact and, when applied correctly, consistently brings near-misses to IoU~$= 1.000$.

\begin{figure*}[t]
  \centering
  \includegraphics[width=\linewidth]{fig3_skills.png}
  \caption{\textbf{Why Skill-Based Composition Instead of a Rigid Workflow.}
  (\textit{Left}) A rigid workflow---a single linear chain---fails when any step encounters a novel failure mode, with no recovery path for heterogeneous cases.
  (\textit{Right}) Our skill-based system: six independent skills (generation, execution, verification, repair, manual check, assembly) all interact with a central verified store. Skills are reusable, replaceable, and independently upgradeable. Hard-tail cases are routed to the manual check skill without disrupting the automated path.}
  \label{fig:skills}
\end{figure*}
